\section{带互补遍历的树码}
\label{sec:tree}

我们沿用 \Cref{sec:preliminaries} 的仿射标记空间 $\mathcal A$ 与行空间
$H=\Aff(\mathcal A,\F)$。
下述编码把仿射行与布尔门表变成一个以标记为指标的词。它的两种互补
遍历控制了在单一标记处取词值的割所看到的相关性。

\begin{definition}[树码]
\label{def:tree-code}
固定整数 $s\ge2$ 与 $m\ge1$，并令
\[
 k=(2s)^m,\qquad d=s^m.
\]
设 $\mathcal T$ 是深度为 $m$ 的完全有根 $2s$ 叉树。
每个内部节点 $v$ 的子节点由 $(g,b)\in[s]\times\{0,1\}$ 索引。其叶子
由 $[k]$ 索引。给每个内部节点指派一张表 $F_v:\F^s\to\{-1,1\}$，
给每个叶子指派一行 $R_i\in H$。
词 $W(R,F)\in\{-1,1\}^{\mathcal A}$ 逐点定义如下。在 $a\in\mathcal A$
处，叶子 $i$ 的符号为 $(-1)^{R_i(a)}$。子符号为 $(Z_{g,b})$ 的内部
节点输出
\begin{equation}
 \left(\prod_{g=1}^s Z_{g,0}\right)F_v(z),
 \qquad (-1)^{z_g}=Z_{g,0}Z_{g,1}.
 \label{eq:gate-rule}
\end{equation}
根节点的输出即 $W(R,F)(a)$。
\end{definition}

我们让这些表被树的第二次遍历所噪声化。
\emph{随机切片}的取法是：先从 $\{0,\ldots,m-1\}$ 中均匀选取 $J_0$。
从根出发，在深度小于 $J_0$ 的每个到达节点处，均匀且独立地选取其
$s$ 对子节点中的一对，并沿该对的两个子节点继续。在深度 $J_0$ 处
到达的节点集 $\mathcal D$ 即为切片。

在 $J_0=j$ 的条件下，切片有 $2^j$ 个节点，而每个固定的深度 $j$ 节点
以概率 $s^{-j}$ 属于它。记 $F_{\rm out}=(F_u)_{u\notin\mathcal D}$ 为
切片外的表。

下文用到的树的两种遍历互为补充。切片构造选取一对并保留两个子节点；
门的 Walsh 展开从每对中保留一个子节点。因此二者的交沿一条单一路径。
Figure~\ref{fig:traversals} 展示了局部规则；逐层迭代即在每个切片深度
得到一个公共节点。

\begin{figure}[ht]
\centering
\begin{tikzpicture}[x=1cm,y=1cm,every node/.style={font=\small},
  mono/.style={draw=orange!80!black,line width=1.4pt},
  slice/.style={draw=blue!65!black,line width=1.4pt,dashed}]
\node[circle,draw,minimum size=7mm] (v) at (0,1.7) {$v$};
\foreach \x/\name/\lab in {-3/a/{(1,0)},-2/b/{(1,1)},-.5/c/{(2,0)},.5/d/{(2,1)},2/e/{(3,0)},3/f/{(3,1)}} {
 \node[circle,draw,fill=white,minimum size=3mm,inner sep=1pt] (\name) at (\x,0) {};
 \node[below] at (\x,-.12) {$\lab$};
 \draw[gray!40] (v) -- (\name);
}
\draw[mono] (v) -- (a);
\draw[mono] (v) -- (d);
\draw[mono] (v) -- (e);
\draw[slice] (v) -- (c);
\draw[slice] (v) -- (d);
\node[circle,draw,fill=violet!65,minimum size=4mm,inner sep=1pt] at (.5,0) {};
\draw[mono] (-3,-1) -- (-2.4,-1);
\node[anchor=west] at (-2.3,-1) {每对一个子节点};
\draw[slice] (-3,-1.5) -- (-2.4,-1.5);
\node[anchor=west] at (-2.3,-1.5) {一对的两个子节点};
\end{tikzpicture}

\caption{一个门处的互补选择，以 $s=3$ 为例。单项式遍历（实线）从每对
中选一个子节点；切片探索（虚线）选一对的两个子节点。恰有一个子节点
被共享，以紫色标出。图中只描绘这一局部选择；同一规则在所选切片之上
逐层重复。}
\label{fig:traversals}
\end{figure}

\begin{lemma}[互补遍历与相关性]
\label{lem:tree-correlation}
对每组固定的表 $F$，根输出作为 $k$ 个叶符号的多重线性多项式只含
次数为 $d$ 的单项式。每个这样的单项式在其到达的每个节点处都从每对
中选一个子节点。每条如此得到的遍历与每个可能的切片 $\mathcal D$
恰交于一个节点。因此，把某一切片上的所有表取反即把 $W$ 取反。

设 $0<t<1$ 与 $0<r<1$。均匀抽取 $F$，并如下构造 $F'$：切片外的表
保持不变，而 $\mathcal D$ 上每对相应的表项耦合为相关为 $t$ 的均匀
符号，各表项间独立。独立地均匀抽取 $R\in H^k$，并以概率 $r$ 保留
$R$ 的每行、否则代之以独立的均匀行（在两个副本中独立进行），得到
$R^1,R^2$。则对每个 $a\in\mathcal A$ 与每个固定的切片，
\begin{equation}
 \E\bigl[W(R^1,F)(a)W(R^2,F')(a)\mid\mathcal D\bigr]
       =t r^{2d}.
 \label{eq:tree-correlation}
\end{equation}
特别地，若 $r=1-\mu/d$（$0<\mu<1/2$），则 $r^{2d}\ge1-2\mu$。
\end{lemma}

\begin{proof}
对固定的门表，按自变量 $z$ 展开：
$F_v(z)=\sum_{I\subseteq[s]}\widehat F_v(I)(-1)^{\sum_{g\in I}z_g}$。
代入 Equation~\eqref{eq:gate-rule}，指标为 $I$ 的项是
\[
 \widehat F_v(I)
 \prod_{g\notin I}Z_{g,0}\prod_{g\in I}Z_{g,1}.
\]
迭代此展开给出
\begin{equation}
 W_F(x)=\sum_{P\in\mathscr P}a_P(F)\prod_{i\in P}x_i,
 \qquad |P|=d,
 \label{eq:tree-polynomial}
\end{equation}
其中 $\mathscr P$ 是在每个到达的内部节点处从每对中选一个子节点所得
叶子集的族。叶子集决定所有这些选择：被选的子节点与该叶子集之交
非空，未选的则不相交。因此不同的遍历给出不同的单项式，而
$a_P(F)$ 是遍历到达的每个节点处 $F_v$ 的一个 Walsh 系数之积。

在根处，单项式遍历与切片探索同时存在。在切片深度之上的任何公共
节点处，两种探索恰选同一个子节点。归纳可得在切片深度处恰有一个
公共节点。于是每个 $a_P$ 恰含来自 $\mathcal D$ 上某张表的一个因子。
把所有这样的表取反，就把 Equation~\eqref{eq:tree-polynomial} 中的
每个系数取反，从而把词取反。此外，切片表的每个 Walsh 系数满足
\[
 \E[\widehat F'_v(I)\mid F_v]=t\widehat F_v(I).
\]
各表之间的条件独立性给出
$\E[a_P(F')\mid F,\mathcal D]=t a_P(F)$。

对固定的 $a$，符号 $x_i^b=(-1)^{R_i^b(a)}$ 均匀且关于 $i$ 独立，
且 $\E[x_i^1x_i^2]=r^2$。于是
Equation~\eqref{eq:tree-polynomial} 中两个单项式乘积的期望在其叶子
集不一致时为零，一致时为 $r^{2d}$。由于 $W_F$ 是布尔值，Parseval
给出 $\sum_P a_P(F)^2=1$。先对行取平均、再对 $F'$ 取平均，即证
Equation~\eqref{eq:tree-correlation}。最后一断言是 Bernoulli
不等式：$(1-\mu/d)^{2d}\ge1-2\mu$。
\end{proof}

该相关性恒等式描述的是在单一标记处取词值的割。为分析任意割，我们
还需要能依赖整张词、却看不到其生成行与表的函数。若
\begin{equation}
 f(R,F)=\phi(W(R,F))
 \label{eq:function-of-code}
\end{equation}
对某个固定的词上函数 $\phi$ 成立，则称 $f$ 为\emph{码的函数}。
这一类在先固定 $\phi$ 再取平均下封闭，包括先把词经固定仿射映射拉回
的取法。对 $r\in[0,1]$，记 $T_r$ 为保留概率 $r$ 的独立行重采样，
并写 $\bar f=T_rf$。该算子保持 $f$ 的值域界。
若 $\phi$ 是奇函数，Lemma~\ref{lem:tree-correlation} 表明 $f$ 与
$\bar f$ 在切片表同时取反下都是奇的。特别地，当所有行与切片外的表
固定时，它们在这些表项上均值为零。

在这些固定数据下，其余随机输入是 $|\mathcal D|2^s$ 个独立表项构成
的立方，由 $(v,z)\in\mathcal D\times\F^s$ 索引。此立方上的傅里叶
次数与影响指\emph{表项}的子集。这与相关性证明中固定门在其 $s$ 比特
自变量 $z$ 上的 Walsh 展开不同。
